\PassOptionsToPackage{unicode}{hyperref}
\PassOptionsToPackage{hyphens}{url}
\PassOptionsToPackage{dvipsnames,svgnames,x11names}{xcolor}
\documentclass[10pt,letterpaper]{article}

\usepackage{xcolor}
\usepackage[margin=2.0cm]{geometry}
\usepackage{lmodern}
\usepackage{amsmath,amssymb,amsfonts,mathtools}
\usepackage{textcomp}
\usepackage{microtype}
\usepackage{parskip}
\usepackage{longtable,booktabs,array,calc}
\usepackage{caption}
\captionsetup[table]{skip=6pt,font=small,labelfont=bf}
\captionsetup[figure]{font=small,labelfont=bf}
\newcounter{none}
\usepackage{etoolbox}
\usepackage{graphicx}
\usepackage{fancyvrb}
\newcommand{\VerbBar}{|}
\newcommand{\VERB}{\Verb[commandchars=\\\{\}]}
\DefineVerbatimEnvironment{Highlighting}{Verbatim}{commandchars=\\\{\},fontsize=\footnotesize}
\RecustomVerbatimEnvironment{verbatim}{Verbatim}{fontsize=\footnotesize}
\newenvironment{Shaded}{}{}
\newcommand{\AlertTok}[1]{\textcolor[rgb]{0.80,0.00,0.00}{\textbf{#1}}}
\newcommand{\AnnotationTok}[1]{\textcolor[rgb]{0.38,0.63,0.69}{\textbf{\textit{#1}}}}
\newcommand{\AttributeTok}[1]{\textcolor[rgb]{0.49,0.56,0.16}{#1}}
\newcommand{\BaseNTok}[1]{\textcolor[rgb]{0.25,0.63,0.44}{#1}}
\newcommand{\BuiltInTok}[1]{\textcolor[rgb]{0.00,0.50,0.00}{#1}}
\newcommand{\CharTok}[1]{\textcolor[rgb]{0.25,0.44,0.63}{#1}}
\newcommand{\CommentTok}[1]{\textcolor[rgb]{0.38,0.63,0.69}{\textit{#1}}}
\newcommand{\CommentVarTok}[1]{\textcolor[rgb]{0.38,0.63,0.69}{\textbf{\textit{#1}}}}
\newcommand{\ConstantTok}[1]{\textcolor[rgb]{0.53,0.00,0.00}{#1}}
\newcommand{\ControlFlowTok}[1]{\textcolor[rgb]{0.00,0.44,0.13}{\textbf{#1}}}
\newcommand{\DataTypeTok}[1]{\textcolor[rgb]{0.56,0.13,0.00}{#1}}
\newcommand{\DecValTok}[1]{\textcolor[rgb]{0.25,0.63,0.44}{#1}}
\newcommand{\DocumentationTok}[1]{\textcolor[rgb]{0.73,0.13,0.13}{\textit{#1}}}
\newcommand{\ErrorTok}[1]{\textcolor[rgb]{0.80,0.00,0.00}{\textbf{#1}}}
\newcommand{\ExtensionTok}[1]{#1}
\newcommand{\FloatTok}[1]{\textcolor[rgb]{0.25,0.63,0.44}{#1}}
\newcommand{\FunctionTok}[1]{\textcolor[rgb]{0.02,0.16,0.49}{#1}}
\newcommand{\ImportTok}[1]{\textcolor[rgb]{0.00,0.50,0.00}{\textbf{#1}}}
\newcommand{\InformationTok}[1]{\textcolor[rgb]{0.38,0.63,0.69}{\textbf{\textit{#1}}}}
\newcommand{\KeywordTok}[1]{\textcolor[rgb]{0.00,0.44,0.13}{\textbf{#1}}}
\newcommand{\NormalTok}[1]{#1}
\newcommand{\OperatorTok}[1]{\textcolor[rgb]{0.40,0.40,0.40}{#1}}
\newcommand{\OtherTok}[1]{\textcolor[rgb]{0.00,0.44,0.13}{#1}}
\newcommand{\PreprocessorTok}[1]{\textcolor[rgb]{0.74,0.48,0.00}{#1}}
\newcommand{\RegionMarkerTok}[1]{#1}
\newcommand{\SpecialCharTok}[1]{\textcolor[rgb]{0.25,0.44,0.63}{#1}}
\newcommand{\SpecialStringTok}[1]{\textcolor[rgb]{0.73,0.40,0.53}{#1}}
\newcommand{\StringTok}[1]{\textcolor[rgb]{0.25,0.44,0.63}{#1}}
\newcommand{\VariableTok}[1]{\textcolor[rgb]{0.10,0.09,0.49}{#1}}
\newcommand{\VerbatimStringTok}[1]{\textcolor[rgb]{0.25,0.44,0.63}{#1}}
\newcommand{\WarningTok}[1]{\textcolor[rgb]{0.38,0.63,0.69}{\textbf{\textit{#1}}}}

\usepackage{newunicodechar}
\newunicodechar{µ}{\ifmmode\mu\else\textmu\fi}
\newunicodechar{°}{\ifmmode^\circ\else\textdegree\fi}
\newunicodechar{·}{\ifmmode\cdot\else\textperiodcentered\fi}
\newunicodechar{²}{\ensuremath{^{2}}}
\newunicodechar{³}{\ensuremath{^{3}}}
\newunicodechar{¹}{\ensuremath{^{1}}}
\newunicodechar{⁴}{\ensuremath{^{4}}}
\newunicodechar{⁶}{\ensuremath{^{6}}}
\newunicodechar{⁻}{\ensuremath{^{-}}}
\newunicodechar{₀}{\ensuremath{_{0}}}
\newunicodechar{ₙ}{\ensuremath{_{n}}}
\newunicodechar{Δ}{\ensuremath{\Delta}}
\newunicodechar{Ω}{\ensuremath{\Omega}}
\newunicodechar{α}{\ensuremath{\alpha}}
\newunicodechar{β}{\ensuremath{\beta}}
\newunicodechar{δ}{\ensuremath{\delta}}
\newunicodechar{ε}{\ensuremath{\varepsilon}}
\newunicodechar{θ}{\ensuremath{\theta}}
\newunicodechar{κ}{\ensuremath{\kappa}}
\newunicodechar{λ}{\ensuremath{\lambda}}
\newunicodechar{π}{\ensuremath{\pi}}
\newunicodechar{ρ}{\ensuremath{\rho}}
\newunicodechar{σ}{\ensuremath{\sigma}}
\newunicodechar{τ}{\ensuremath{\tau}}
\newunicodechar{→}{\ensuremath{\rightarrow}}
\newunicodechar{∎}{\ensuremath{\blacksquare}}
\newunicodechar{−}{\ensuremath{-}}
\newunicodechar{∧}{\ensuremath{\wedge}}
\newunicodechar{∫}{\ensuremath{\int}}
\newunicodechar{≈}{\ensuremath{\approx}}
\newunicodechar{≤}{\ensuremath{\leq}}
\newunicodechar{≥}{\ensuremath{\geq}}
\newunicodechar{⊆}{\ensuremath{\subseteq}}
\newunicodechar{✓}{\ensuremath{\checkmark}}
\newunicodechar{≠}{\ensuremath{\neq}}
\newunicodechar{↔}{\ensuremath{\leftrightarrow}}
\newunicodechar{≪}{\ensuremath{\ll}}
\newunicodechar{↛}{\ensuremath{\nrightarrow}}
\newunicodechar{⇒}{\ensuremath{\Rightarrow}}
\newunicodechar{∪}{\ensuremath{\cup}}
\newunicodechar{⟨}{\ensuremath{\langle}}
\newunicodechar{⟩}{\ensuremath{\rangle}}
\newunicodechar{・}{\textperiodcentered}

\usepackage{xurl}
\usepackage{bookmark}
\usepackage{hyperref}

\setcounter{secnumdepth}{-\maxdimen}
\providecommand{\tightlist}{\setlength{\itemsep}{1pt}\setlength{\parskip}{0pt}}

\definecolor{openieAmber}{HTML}{D97706}
\definecolor{openieAmberBg}{HTML}{FFFBEB}
\definecolor{openieAmberBorder}{HTML}{F59E0B}
\definecolor{openieSlate}{HTML}{334155}
\definecolor{openieDark}{HTML}{0F172A}
\definecolor{openieCodeBg}{HTML}{F8FAFC}
\definecolor{openieCodeBorder}{HTML}{CBD5E1}
\definecolor{openieBlue}{HTML}{0057B8}

\hypersetup{
  pdftitle={"AI" Reads Human: Human Language Encoding at Machine Speed and the Energy Cost},
  pdfauthor={David Charlot, Open Interface Engineering},
  colorlinks=true,
  linkcolor=openieBlue,
  citecolor=openieBlue,
  urlcolor=openieBlue,
  pdfcreator={XeTeX / Tectonic — Open Interface Engineering}
}

\urlstyle{same}
\setlength{\emergencystretch}{3em}
\hfuzz=6pt
\sloppy

\newcommand{\openiefig}[3]{%
  \IfFileExists{#1}{%
    \includegraphics[width=#2\linewidth,height=0.70\textheight,keepaspectratio]{#1}%
  }{%
    \fcolorbox{openieCodeBorder}{openieCodeBg}{%
      \parbox{0.92\linewidth}{\footnotesize\ttfamily Figure asset embedded in archival PDF; interactive/vector source: \url{#3}}%
    }%
  }%
}

\title{\vspace{-1.2em}\textbf{"AI" Reads Human: Human Language Encoding at Machine Speed and the Energy Cost}}
\author{\textbf{David Charlot, Open Interface Engineering}\\[0.15em]
{\small Work in progress $\cdot$ Open Interface Engineering $\cdot$ \href{https://research.openie.dev}{research.openie.dev}}}
\date{\small Revision \texttt{af57f48} (22 November 2026) $\cdot$ Web: \href{https://research.openie.dev/papers/ai-reads-human/}{research.openie.dev/papers/ai-reads-human/} $\cdot$ Archival LaTeX: \href{https://research.openie.dev/pdfs/latex/ai-reads-human.pdf}{/pdfs/latex/ai-reads-human.pdf} (\href{https://research.openie.dev/pdfs/latex/ai-reads-human.tex}{.tex})}

\begin{document}
\maketitle
\vspace{-0.6em}

% OpenIE Measurement Honesty Callout Box
\noindent\fcolorbox{openieAmberBorder}{openieAmberBg}{%
  \parbox{\dimexpr\linewidth-2\fboxsep-2\fboxrule\relax}{%
    \small\color{openieDark}%
    \textbf{Measurement.} Research study, not a journal final. Parse time, token counts, and Apple M5 Max dual-channel \texttt{reported\_j} (\texttt{L3 proc\_pid\_rusage(RUSAGE\_INFO\_V6).ri\_energy\_nj} and \texttt{L2 macmon IOReport} across \textbf{\texttt{34} workloads} and \texttt{170} windows on the reproducible \texttt{20,000}-record corpus: \textbf{\texttt{280.4 nJ/record}} fixed binary \texttt{struct} to \textbf{\texttt{2,212.1 nJ/record}} stdlib \texttt{JSON}, \textbf{\texttt{69,922.1 nJ/record}} C \texttt{libyaml}, \textbf{\texttt{575,799.3 nJ/record}} pure-Python \texttt{YAML}, and \textbf{\texttt{27,678.7 nJ/record}} \texttt{cl100k\_base} \texttt{JSON} tokenization; \texttt{Falsifier 6} triggered on \texttt{109/170} shared-host \texttt{L2} windows) were measured locally. No hardware power rail or wall meter was attached (\texttt{measured\_j = None}, \texttt{board\_metered = false}); model energy figures are other authors' measurements.
  }%
}

\vspace{6pt}
% OpenIE Amber Left-Rule Summary Quote
\noindent{\color{openieAmber}\vrule width 3pt}\hspace{8pt}%
\parbox{\dimexpr\linewidth-11pt\relax}{%
  \small\itshape\color{openieSlate}%
  Work in progress. Why so much of what machines exchange is written for people, what that costs in bytes, tokens, time, and per-process Apple Silicon kernel rusage energy (reported\_j) across 34 parse and tokenize workloads, and why external wall-meter energy (measured\_j) remains open.
}
\vspace{6pt}

{
\small
\setcounter{tocdepth}{2}
\tableofcontents
}
\vspace{4pt}
\hrule
\vspace{8pt}

\emph{Work in progress, revision
\href{https://github.com/openIE-dev/research-openie-web/commits/main/src/content/papers/ai-reads-human.md}{\texttt{af57f48}},
22 November 2026. Interactive observatory, data, code and read ledger:
\href{https://research.openie.dev/living/ai-reads-human/}{research.openie.dev/living/ai-reads-human/}.
Shared terms: \href{https://research.openie.dev/glossary/}{glossary}.}

\section{Abstract}\label{abstract}

Software agents pass data to each other in formats made for people to
read, such as JSON, YAML, Markdown, XML and CSV. This study names the
extra work this causes the impedance tax, defines it with units, and
measures three parts of it on a seeded 20,000-record corpus: parse time,
token count across five open tokenizers, and dual-channel Apple M5 Max
operating-system/on-chip \texttt{reported\_j} counters across all 34
parse and \texttt{tiktoken} workloads (section 5c). While an initial
unisolated whole-chip attempt on two busy laptops could not separate
parsing from background co-tenant drift (section 5b), isolating each
workload inside a dedicated worker subprocess with macOS kernel
per-process \texttt{proc\_\allowbreak{}pid\_\allowbreak{}rusage(\allowbreak{}RUSAGE\_\allowbreak{}INFO\_\allowbreak{}V6).\allowbreak{}ri\_\allowbreak{}energy\_\allowbreak{}nj}
(\texttt{L3}, label \texttt{D}) resolves a \texttt{2,053.8×}
parse-energy span --- from \texttt{280.\allowbreak{}4\ nJ/\allowbreak{}record} for fixed binary
(\texttt{struct.\allowbreak{}iter\_\allowbreak{}unpack}) to \texttt{2,\allowbreak{}212.\allowbreak{}1\ nJ/\allowbreak{}record} for stdlib
JSON (\texttt{7.89×}), \texttt{69,\allowbreak{}922.\allowbreak{}1\ nJ/\allowbreak{}record} for C
\texttt{libyaml} (\texttt{249.4×}), and \texttt{575,\allowbreak{}799.\allowbreak{}3\ nJ/\allowbreak{}record}
for pure-Python YAML (\texttt{2,053.8×}) --- and shows that
\texttt{cl100k\_base} tokenization of compact JSON
(\texttt{27,\allowbreak{}679.\allowbreak{}2\ nJ/\allowbreak{}record}, \texttt{98.7×} fixed binary) spends
\texttt{12.5×} the on-chip \texttt{reported\_j} of \texttt{json.loads()}
itself. Concurrently, whole-chip \texttt{IOReport} (\texttt{L2}, label
\texttt{A2}) triggered Falsifier 6 (\texttt{\textgreater{}\ 5.0\%}
pre/post idle baseline drift) on \texttt{109} of \texttt{170}
shared-host measurement windows, confirming why unisolated package
subtraction fails on active hosts. No external inline power probe or
wall meter was attached (\texttt{measured\_\allowbreak{}j\ =\ None},
\texttt{board\_\allowbreak{}metered\ =\ false}), and GPU model prefill/decode joules
remain planned (M4).

\textbf{In plain words}

\begin{enumerate}
\def\labelenumi{\arabic{enumi}.}
\tightlist
\item
  Computers and people often use the same written format. That format is
  easy for people to read, but the computer has to do extra work to read
  it back.
\item
  On our 20,000-record test corpus, reading fixed binary
  (\texttt{struct.\allowbreak{}iter\_\allowbreak{}unpack}) cost \texttt{280.\allowbreak{}4\ nJ/\allowbreak{}record} in Apple
  M5 Max per-process kernel \texttt{ri\_energy\_nj}
  (\texttt{reported\_j}), compared with \texttt{2,\allowbreak{}212\ nJ/\allowbreak{}record}
  (\texttt{7.89×}) for Python's built-in JSON reader,
  \texttt{69,\allowbreak{}922\ nJ/\allowbreak{}record} (\texttt{249.4×}) for C-accelerated YAML,
  and \texttt{575,\allowbreak{}799\ nJ/\allowbreak{}record} (\texttt{2,053.8×}) for pure-Python
  YAML.
\item
  Tokenizing compact JSON for an AI model with \texttt{tiktoken}
  (\texttt{cl100k\_base}) cost \texttt{27,\allowbreak{}679\ nJ/\allowbreak{}record}
  (\texttt{1.69×} CSV tokenization and \texttt{12.5×}
  \texttt{json.loads()} parsing itself), before the language model runs
  a single matrix multiply. All energy figures on this page are
  kernel/on-chip \texttt{reported\_j} counters
  (\texttt{measured\_\allowbreak{}j\ =\ None}, \texttt{board\_\allowbreak{}metered\ =\ false});
  GPU model prefill/decode energy and external hardware power probes are
  still planned.
\end{enumerate}

\begin{quote}
\textbf{Status: work in progress, and it will grow.} Sections marked
\emph{Planned} (such as GPU prefill/decode counters M4, natural-language
corpora M6, and cross-language C/Rust parsers M8) are not done. Every
number on this page is recomputed from downloadable files under
\texttt{/\allowbreak{}living/\allowbreak{}ai-\allowbreak{}reads-\allowbreak{}human/\allowbreak{}}. On-chip energy counters are strictly
\texttt{reported\_j} (\texttt{measured\_\allowbreak{}j\ =\ None},
\texttt{board\_\allowbreak{}metered\ =\ false}).
\end{quote}

\section{1. Background: formats written for
people}\label{background-formats-written-for-people}

Most data that agents exchange today travels in formats whose
specifications put people ahead of machines. The passages cited below
were read for this page; each source's read status (FULL, PARTIAL or
UNAVAILABLE) is in the
\href{https://research.openie.dev/living/ai-reads-human/READ_LEDGER.md}{read
ledger}.

\begin{itemize}
\tightlist
\item
  \textbf{TeX} began in 1982: the source file \texttt{tex.web} (header,
  first 19,000 bytes read; PARTIAL) records that version 0 was released
  in September 1982 and dates later versions from 3.0 (March 1990)
  onward.
\item
  \textbf{Markdown} was developed by John Gruber with help from Aaron
  Swartz and released in 2004 as a syntax description and a Perl script
  that converts to HTML; the CommonMark specification (version 0.31.2,
  28 January 2024) says that description ``does not specify the syntax
  unambiguously'' and lists questions it leaves open (CommonMark 0.31.2,
  sections 1.1 and 1.2).
\item
  \textbf{CSV} was documented in October 2005 by RFC 4180, an
  \emph{informational} RFC that says ``there is no formal specification
  in existence'' and writes down what most implementations seem to do.
\item
  \textbf{JSON} is defined by RFC 8259 (December 2017, obsoleting RFC
  7159) as ``a lightweight, text-based, language-independent data
  interchange format.'' RFC 7493 (March 2015, I-JSON) narrows it for
  protocols.
\item
  \textbf{Canonical and binary forms} also exist in the standards
  record: RFC 8785 (June 2020) defines a JSON Canonicalization Scheme,
  and RFC 8949 (December 2020, obsoleting RFC 7049) defines CBOR, a
  binary data format with a section on deterministic encoding (section
  4.2) that protocols may adopt.
\item
  \textbf{YAML} revision 1.2.2 (1 October 2021) lists its design goals
  ``in decreasing priority,'' and the highest-priority goal is ``YAML
  should be easily readable by humans'' (section 1.1).
\end{itemize}


\begin{center}
\openiefig{fig/ai-reads-human-fig01-format-timeline.png}{0.94}{https://research.openie.dev/living/ai-reads-human/charts/format-timeline.svg}
\end{center}


\emph{Figure 1. Dates printed in the documents listed in the read
ledger. Generated by \texttt{scripts/\allowbreak{}make\_\allowbreak{}charts.\allowbreak{}py}.}

These facts do not say that human-readable formats are mistakes. They
say that readability by people was a stated design goal, and that the
cost of reading them back by machine was not the thing being minimised.
Where other work reaches machine-oriented forms (CBOR, canonical JSON,
Protocol Buffers, Avro, safetensors, GGUF), this study records it as a
convergence observation: several communities arrived at fixed or
deterministic encodings for machine exchange while keeping text for
people.

\section{2. The impedance tax, defined}\label{the-impedance-tax-defined}

Let a set of records R carry information that a machine consumer needs.
Let F be a format with encoder enc\_F and decoder dec\_F, and let M be a
canonical machine form: one valid byte string per value.

{\def\LTcaptype{none} % do not increment counter
{\footnotesize\setlength{\tabcolsep}{3.8pt}\renewcommand{\arraystretch}{1.15}
\begin{longtable}[]{@{}
  >{\raggedright\arraybackslash}p{(\linewidth - 4\tabcolsep) * \real{0.2268}}
  >{\raggedright\arraybackslash}p{(\linewidth - 4\tabcolsep) * \real{0.5670}}
  >{\raggedright\arraybackslash}p{(\linewidth - 4\tabcolsep) * \real{0.2062}}@{}}
\toprule\noalign{}
\begin{minipage}[b]{\linewidth}\raggedright
\textbf{Symbol}
\end{minipage} & \begin{minipage}[b]{\linewidth}\raggedright
\textbf{Meaning}
\end{minipage} & \begin{minipage}[b]{\linewidth}\raggedright
\textbf{Unit}
\end{minipage} \\
\midrule\noalign{}
\endhead
\bottomrule\noalign{}
\endlastfoot
B\_F & encoded size of R in format F & byte \\
N & number of records & 1 \\
k & a tokenizer & n/a \\
T\_\{F,k\} & tokens produced by k on the encoding of R & token \\
d\_\{F,k\} = B\_F / T\_\{F,k\} & bytes per token & byte/token \\
t\_parse(F) & wall time to decode R to in-memory values & s \\
P & average power during an operation & W = J/s \\
E & energy of an operation, E = ∫ P dt, approximately P\_avg · t & J \\
e\_prefill & model energy per input token & J/token \\
\end{longtable}
}
}

The \textbf{impedance tax} of F relative to M is the extra energy a
consumer spends because the information arrives in F:

Z(F) = E\_enc(F) + E\_parse(F) + E\_tok(F,k) + e\_prefill · T\_\{F,k\} +
E\_ctx(F)

τ(F) = Z(F) − Z(M) (joules), θ(F) = Z(F) / Z(M) (dimensionless)

E\_ctx(F) is the extra decode cost that a longer context causes; its
coefficient depends on the model and is a measurement target (M4 below).
Section 5c measures \texttt{E\_parse(F)} and \texttt{E\_tok(F,k)} in
per-process kernel \texttt{reported\_j} (\texttt{ri\_energy\_nj}) on
Apple M5 Max.

Three dimensionless ratios can be measured without any energy counter:

\begin{itemize}
\tightlist
\item
  byte ratio β\_F = B\_F / B\_ref
\item
  token ratio κ\_\{F,k\} = T\_\{F,k\} / T\_\{ref,k\}
\item
  parse-time ratio π\_F = t\_parse(F) / t\_parse(ref)
\end{itemize}

and they are linked by one identity that the token numbers below obey:

κ\_F = β\_F · (d\_ref / d\_F)

so a format can cost more tokens because it has more bytes, because its
bytes split into more pieces, or both.


\begin{center}
\openiefig{fig/ai-reads-human-fig02-impedance-pipeline.png}{0.94}{https://research.openie.dev/living/ai-reads-human/charts/impedance-pipeline.svg}
\end{center}


\emph{Figure 2. The terms of Z(F). With no model in the path (machine to
machine) only encode, transport and parse apply. With a model reading
text, tokenize and read apply.}

\subsection{2.1 Physical scale, and what this page does not
claim}\label{physical-scale-and-what-this-page-does-not-claim}

The thermodynamic floor for erasing one bit at temperature T is the
Landauer limit E = k\_B T ln 2. With k\_B = 1.380649e-23 J/K (exact by
the 2019 SI definition; the value was checked on 18 October 2026 against
the NIST CODATA 2022 page for the Boltzmann constant) and T = 300 K this
is \textbf{2.871e-21 J} per bit, recomputed in
\texttt{scripts/\allowbreak{}verify\_\allowbreak{}numbers.\allowbreak{}py}. Eight bits is 2.297e-20 J. This is
a lower bound for irreversible logic, not an estimate of what parsing
costs: real software pays orders of magnitude more, in memory traffic,
instruction execution and cooling. This page does not use the limit to
predict any number.

Energy figures for model inference come from other authors' measurements
on their hardware, and are background only, not our measurement:

\begin{itemize}
\tightlist
\item
  Samsi et al., ``From Words to Watts: Benchmarking the Energy Costs of
  Large Language Model Inference'' (arXiv:2310.03003), define energy per
  output token as total GPU energy divided by decoded output tokens, and
  report about 3 to 4 J per output token for a 65-billion-parameter
  LLaMA model at a maximum generation length of 512, measured with
  NVIDIA DCGM on their MIT SuperCloud setup (V100 GPUs for the 8, 16 and
  32 shard runs). That paper does not separate prefill from decode
  energy.
\item
  Luccioni, Jernite and Strubell, ``Power Hungry Processing: Watts
  Driving the Cost of AI Deployment?'' (ACM FAccT 2024, DOI
  10.1145/3630106.3658542), measured energy per 1,000 inferences by task
  with CodeCarbon on an 8-GPU A100 node; their Table 2 reports mean
  0.002 kWh for text classification and 0.047 kWh for text generation.
  These numbers were checked against the ACM Corrected Version of Record
  V1.1 (6 November 2024), which gives the same Table 2 values and the
  same setup (8 A100 80GB GPUs on AWS, CodeCarbon). Cite V1.1: the paper
  notes that an earlier version used 0.012 kWh in one comparison. The
  rest of the paper was read only in part.
\end{itemize}

A per-word DRAM-read energy figure (a widely quoted 640 pJ per 32-bit
read) is \textbf{not used on this page}. Its attributed source, Horowitz
(2014), was read: the paper gives about 1 to 2 nJ per DRAM access and
the plenary slides 1.3 to 2.6 nJ per 64-bit access; the 640 pJ figure
appears only in later slides that cite it. Section 9 records this
status.

\section{3. What was measured}\label{what-was-measured}

Initial parse-time and token-count baselines ran on one virtualised
Intel Xeon machine (8 vCPU, Linux 6.12, Python 3.13.5), which exposes no
RAPL energy counters (\texttt{/\allowbreak{}sys/\allowbreak{}class/\allowbreak{}powercap} is absent), followed
by Apple M5 Max energy campaigns reported in sections 5b and 5c. The
corpus is synthetic and seeded (\texttt{scripts/\allowbreak{}gen\_\allowbreak{}corpus.\allowbreak{}py}, seed
20261016): 20,000 records, each with an id, a three-word name, a float
score, a boolean, three tags and a 31-bit integer.


\begin{center}
\openiefig{fig/ai-reads-human-fig03-one-record-five-ways.png}{0.94}{https://research.openie.dev/living/ai-reads-human/charts/one-record-five-ways.svg}
\end{center}


\emph{Figure 3. Record 0 in CSV, compact JSON, YAML, XML and a 23-byte
fixed binary layout. The binary layout stores each word as a one-byte
index into a shared list of 24 words, so reading it back needs that
list. Byte counts of the text forms include the line break.}

\subsection{3.1 Parse time}\label{parse-time}


\begin{center}
\openiefig{fig/ai-reads-human-fig04-parse-time-headline.png}{0.94}{https://research.openie.dev/living/ai-reads-human/charts/parse-time-headline.svg}
\end{center}


\emph{Figure 4. Minimum of 9 repeats over 20,000 records on the Linux
Xeon baseline. Source:
\texttt{/\allowbreak{}living/\allowbreak{}ai-\allowbreak{}reads-\allowbreak{}human/\allowbreak{}data/\allowbreak{}parse\_\allowbreak{}time.\allowbreak{}tsv}, column
\texttt{ns\_per\_record}. Log scale.}

{\def\LTcaptype{none} % do not increment counter
{\footnotesize\setlength{\tabcolsep}{3.8pt}\renewcommand{\arraystretch}{1.15}
\begin{longtable}[]{@{}
  >{\raggedright\arraybackslash}p{(\linewidth - 4\tabcolsep) * \real{0.3854}}
  >{\raggedright\arraybackslash}p{(\linewidth - 4\tabcolsep) * \real{0.3021}}
  >{\raggedright\arraybackslash}p{(\linewidth - 4\tabcolsep) * \real{0.3125}}@{}}
\toprule\noalign{}
\begin{minipage}[b]{\linewidth}\raggedright
\textbf{Format and parser}
\end{minipage} & \begin{minipage}[b]{\linewidth}\raggedright
\textbf{ns per record (Xeon)}
\end{minipage} & \begin{minipage}[b]{\linewidth}\raggedright
\textbf{ratio to fixed binary}
\end{minipage} \\
\midrule\noalign{}
\endhead
\bottomrule\noalign{}
\endlastfoot
Fixed binary, \texttt{struct.\allowbreak{}iter\_\allowbreak{}unpack} & 116 & 1.0 \\
JSON, Python standard library & 904 & 7.8 \\
YAML, PyYAML with libyaml & 85,300 & 735.3 \\
\end{longtable}
}
}

YAML took 94.4 times as long as stdlib JSON on the same records on the
Xeon host (and \texttt{22,404\ ns} vs \texttt{570\ ns} = \texttt{39.3×}
for C \texttt{libyaml}, or \texttt{161,227\ ns} = \texttt{282.9×} for
pure-Python YAML, on Apple M5 Max in section 5c). The full Xeon table of
every format and parser tried is in Figure 5.


\begin{center}
\openiefig{fig/ai-reads-human-fig05-parse-time-all.png}{0.94}{https://research.openie.dev/living/ai-reads-human/charts/parse-time-all.svg}
\end{center}


\emph{Figure 5. All rows of \texttt{data/\allowbreak{}parse\_\allowbreak{}time.\allowbreak{}tsv} except the
zero-copy numpy row, which builds no objects and is not comparable.}

\textbf{Limits of this measurement.} (a) Parsers do different work: the
fixed binary reader returns tuples, JSON and YAML build typed Python
objects, CSV returns strings, and the XML and Markdown rows build trees
or token streams without converting values, so the ratios compare
end-to-end library behaviour, not the intrinsic cost of each format. (b)
One language runtime (CPython 3.13). (c) The corpus is synthetic and
identifier-light. (d) Time alone is not energy: section 5c therefore
measures per-process kernel \texttt{ri\_energy\_nj}
(\texttt{reported\_j}) directly across all 34 parse and tokenize
workloads.

\subsection{3.2 Tokens}\label{tokens}

A language model reads tokens, not bytes. The same records, counted with
the \texttt{cl100k\_base} tokenizer on the initial 2,000 records of each
view (and verified across all 20,000 records in section 5c):


\begin{center}
\openiefig{fig/ai-reads-human-fig06-tokens-vs-csv.png}{0.94}{https://research.openie.dev/living/ai-reads-human/charts/tokens-vs-csv.svg}
\end{center}


\emph{Figure 6. κ = tokens per record of the view divided by tokens per
record of CSV. Source: \texttt{/\allowbreak{}living/\allowbreak{}ai-\allowbreak{}reads-\allowbreak{}human/\allowbreak{}data/\allowbreak{}tokens.\allowbreak{}tsv}.}

{\def\LTcaptype{none} % do not increment counter
{\footnotesize\setlength{\tabcolsep}{3.8pt}\renewcommand{\arraystretch}{1.15}
\begin{longtable}[]{@{}
  >{\raggedright\arraybackslash}p{(\linewidth - 4\tabcolsep) * \real{0.3958}}
  >{\raggedright\arraybackslash}p{(\linewidth - 4\tabcolsep) * \real{0.2917}}
  >{\raggedright\arraybackslash}p{(\linewidth - 4\tabcolsep) * \real{0.3125}}@{}}
\toprule\noalign{}
\begin{minipage}[b]{\linewidth}\raggedright
\textbf{View}
\end{minipage} & \begin{minipage}[b]{\linewidth}\raggedright
\textbf{tokens per record}
\end{minipage} & \begin{minipage}[b]{\linewidth}\raggedright
\textbf{κ (relative to CSV)}
\end{minipage} \\
\midrule\noalign{}
\endhead
\bottomrule\noalign{}
\endlastfoot
CSV & 22.9 & 1.00 \\
JSON, compact & 36.2 & 1.58 \\
YAML & 51.4 & 2.24 \\
JSON, pretty-printed (indent 2) & 65.6 & 2.86 \\
\end{longtable}
}
}

The identity in section 2 holds on these data: for compact JSON, β =
1.912 and d\_CSV/d\_JSON = 0.826, giving κ = 1.58; for YAML, 2.028 ×
1.107 = 2.24; for pretty JSON, 3.056 × 0.936 = 2.86. Compact JSON costs
more tokens mainly because it has about 1.9 times the bytes of CSV, and
pretty JSON mostly for the same reason plus indentation. YAML has both
more bytes and shorter tokens.


\begin{center}
\openiefig{fig/ai-reads-human-fig07-tokens-five-tokenizers.png}{0.94}{https://research.openie.dev/living/ai-reads-human/charts/tokens-five-tokenizers.svg}
\end{center}


\emph{Figure 7. The same ratio under \texttt{cl100k\_base},
\texttt{o200k\_base}, GPT-2, Mistral-7B-v0.1 and Qwen2.5-0.5B
tokenizers. The ordering holds in all five; the size of the gap does not
(pretty JSON is 3.66 times CSV under GPT-2 and 2.22 times under
Qwen2.5). No closed commercial tokenizer was tested.}


\begin{center}
\openiefig{fig/ai-reads-human-fig08-bytes-per-record.png}{0.94}{https://research.openie.dev/living/ai-reads-human/charts/bytes-per-record.svg}
\end{center}


\emph{Figure 8. Bytes per record. The 23-byte binary figure does not
count the shared word list; any honest size claim for a binary form must
name the list (or dictionary) it depends on, by version and hash.}

\textbf{A negative result worth stating.} Writing the binary layout as
text for a model does not remove the tax: as base64 it was 20.3 tokens
per record (κ = 0.89) and as hex 22.7 (κ = 0.99), about the same as CSV,
because a text model needs a text view. The saving for a model comes
from sending fewer, better chosen fields, not from changing the encoding
of what is sent. That claim is Planned (M7), not measured.

\section{4. Why the hazards come with the
readability}\label{why-the-hazards-come-with-the-readability}

Formats designed for people and for loosely coupled programs let one
input mean different things to different readers. The table lists
hazards found in the read specifications, and in tests run on named
library versions (Python 3.13.5, Node v20.19.2, jq 1.7, PyYAML 6.0.2,
lxml 6.1.3). Test lines (T, U, X, Y, C codes) are in
\texttt{tests-\allowbreak{}formats/\allowbreak{}RESULTS.\allowbreak{}txt} and describe only those versions.
``Spec'' means the text of the cited document.

{\def\LTcaptype{none} % do not increment counter
{\footnotesize\setlength{\tabcolsep}{3.8pt}\renewcommand{\arraystretch}{1.15}
\begin{longtable}[]{@{}
  >{\raggedright\arraybackslash}p{(\linewidth - 6\tabcolsep) * \real{0.1875}}
  >{\raggedright\arraybackslash}p{(\linewidth - 6\tabcolsep) * \real{0.1250}}
  >{\raggedright\arraybackslash}p{(\linewidth - 6\tabcolsep) * \real{0.4688}}
  >{\raggedright\arraybackslash}p{(\linewidth - 6\tabcolsep) * \real{0.2188}}@{}}
\toprule\noalign{}
\begin{minipage}[b]{\linewidth}\raggedright
\textbf{Class}
\end{minipage} & \begin{minipage}[b]{\linewidth}\raggedright
\textbf{Format}
\end{minipage} & \begin{minipage}[b]{\linewidth}\raggedright
\textbf{Hazard}
\end{minipage} & \begin{minipage}[b]{\linewidth}\raggedright
\textbf{Evidence}
\end{minipage} \\
\midrule\noalign{}
\endhead
\bottomrule\noalign{}
\endlastfoot
Ambiguity & JSON & Duplicate member names give ``unpredictable''
behaviour; implementations report the last, all, or fail. Three tools
tested kept the last value. & RFC 8259 section 4; T1 \\
Ambiguity & Markdown & The original description does not specify syntax
unambiguously; dialects diverged. & CommonMark 0.31.2 section 1.2 \\
Parser differential & JSON numbers & 9007199254740993 is exact in Python
and rounds to \ldots992 in Node. & RFC 8259 section 6; T2 and N2 \\
Parser differential & JSON & Python's \texttt{json} accepts bare
\texttt{NaN} and \texttt{Infinity}, which the spec says are not
permitted. & RFC 8259 section 6; T6 \\
Non-canonical forms & JSON & \texttt{1}, \texttt{1.0}, \texttt{1e0},
\texttt{10E-1}, \texttt{0.1e1} are one number written five ways, and two
Python types. & U3 \\
Non-canonical forms & JSON & Member order does not change the meaning
but changes the bytes and any hash of them. & T8 \\
Ordering & JSON canonicalization & Sorting by UTF-16 code units differs
from code point order for U+1F600 versus U+FFEE. & RFC 8785 section
3.2.3; U2 \\
Non-canonical forms & CBOR & Deterministic encoding is an optional
profile that protocols may adopt; it is not the default. & RFC 8949
section 4.2 \\
Non-canonical forms & Protocol Buffers & ``Do not assume the byte output
of a serialized message is stable.'' & Protobuf encoding guide,
``Implications'' \\
Number hazards & JSON & Implementations may limit range and precision;
the exactly representable integer range is ±(2\^{}53)−1. & RFC 8259
sections 6 and 9; RFC 7493 section 2.2 \\
Number hazards & JSON & \texttt{1e400} becomes \texttt{inf} in both
Python and Node; \texttt{-0} is handled differently. & T3, N3, T4, N4 \\
Number hazards & YAML 1.1 resolver & In PyYAML, \texttt{010} is 8,
\texttt{12:30:45} is 45045 and \texttt{country:\ NO} is
\texttt{\{country:\ False\}}. & Y rows \\
Number hazards & TOML & The spec says an integer that cannot be stored
losslessly is an error; Python's \texttt{tomllib} accepted 2\^{}63. &
TOML 1.0.0; T rows \\
Unicode & JSON & \texttt{"\textbackslash{}uDEAD"} (an unpaired
surrogate) is allowed by the grammar; it parses, then fails on UTF-8
encode. & RFC 8259 section 8.2; T5 \\
Unicode & JSON & NFC and NFD forms of the same text are distinct keys. &
U1 \\
Injection & CSV & A cell starting with \texttt{=} is plain data to the
parser and a formula to some spreadsheet programs. & OWASP CSV
Injection; C row \\
Injection & XML & With entity resolution on, lxml read a local file
through an external entity. & X rows \\
Resource blow-up & YAML & 166 bytes of aliases expanded to 59,049
logical leaves. & Y alias test \\
Resource blow-up & XML & 202 bytes of internal entities expanded to
10,000 characters (49.5 times). & X rows \\
Resource blow-up & JSON & Depth limits are implementation-defined;
Python raised \texttt{RecursionError} at 100,000 nesting. & RFC 8259
section 9; T9 \\
Totality & any parser & Library exceptions are control flow outside the
function's return type. & T5, T9 \\
\end{longtable}
}
}

None of these is a reason to avoid the formats for people. Each is a
reason not to use the same text as the single authority between
machines.

\section{5. Planned: measuring energy}\label{planned-measuring-energy}

The measurement roadmap separates completed parse/tokenize
\texttt{reported\_j} counters from open GPU inference and external
hardware power probes:

{\def\LTcaptype{none} % do not increment counter
{\footnotesize\setlength{\tabcolsep}{3.8pt}\renewcommand{\arraystretch}{1.15}
\begin{longtable}[]{@{}
  >{\raggedright\arraybackslash}p{(\linewidth - 4\tabcolsep) * \real{0.0625}}
  >{\raggedright\arraybackslash}p{(\linewidth - 4\tabcolsep) * \real{0.5521}}
  >{\raggedright\arraybackslash}p{(\linewidth - 4\tabcolsep) * \real{0.3854}}@{}}
\toprule\noalign{}
\begin{minipage}[b]{\linewidth}\raggedright
\textbf{ID}
\end{minipage} & \begin{minipage}[b]{\linewidth}\raggedright
\textbf{Measurement}
\end{minipage} & \begin{minipage}[b]{\linewidth}\raggedright
\textbf{Status}
\end{minipage} \\
\midrule\noalign{}
\endhead
\bottomrule\noalign{}
\endlastfoot
M1, M2, M3 & Parse time, tokens, tokenizer throughput & \textbf{Done}
(\texttt{data/\allowbreak{}parse\_\allowbreak{}time.\allowbreak{}tsv}, \texttt{data/tokens.tsv},
\texttt{data/\allowbreak{}tokenize\_\allowbreak{}time.\allowbreak{}tsv}) \\
M5 & Parse and \texttt{tiktoken} energy across all 34 workloads via
dual-channel Apple M5 Max \texttt{reported\_j} (\texttt{L3} per-process
\texttt{ri\_energy\_nj} + \texttt{L2} whole-chip \texttt{IOReport} with
Falsifier 6 idle-drift gating) & \textbf{Done for \texttt{reported\_j}}
(section 5c; \texttt{data/\allowbreak{}energy\_\allowbreak{}reported\_\allowbreak{}j.\allowbreak{}tsv};
\texttt{measured\_\allowbreak{}j\ =\ None}, \texttt{board\_\allowbreak{}metered\ =\ false}) \\
M4 & Prefill energy and decode energy per token against context length,
same model, GPU with a hardware energy counter or external rail/wall
probe & Planned (needs dedicated GPU counter or inline power probe) \\
M6 & A second corpus: natural-language documentation and logs against
identifier-heavy telemetry & Planned \\
M7 & Task level: tokens an agent needs to answer a fixed question set
from a projection of the fields it needs versus the full document &
Planned \\
M8 & Other parsers per format (simdjson, rapidyaml, libxml2 C,
pulldown-cmark) to bound language-runtime effects & Planned \\
\end{longtable}
}
}

External wall-plug or inline USB-C power-probe readings
(\texttt{measured\_j}) remain uncollected; section 5c reports
operating-system/on-chip \texttt{reported\_j} counters strictly.

\section{5b. Measured energy: no net effect
resolved}\label{b.-measured-energy-no-net-effect-resolved}

\emph{First whole-chip subtraction attempt, run on 2026-10-18 on two
Apple M5 Max laptops (48 GB on battery power, 128 GB on AC), macOS
27.0.1, Python 3.13.16, power counter macmon 0.10.0 (48 GB) and 0.9.0
(128 GB), reading whole-SoC \texttt{IOReport} and SMC sensors
(\texttt{all\_power} and \texttt{sys\_power}). Nine parse workloads on
the same 20,000 records, each looped for 10 s on one thread, 10
interleaved rounds, with a 10 s idle window before and after every
round. Files: \texttt{/\allowbreak{}living/\allowbreak{}ai-\allowbreak{}reads-\allowbreak{}human/\allowbreak{}data/\allowbreak{}energy\_\allowbreak{}joules.\allowbreak{}tsv}
and raw samples in
\texttt{/\allowbreak{}living/\allowbreak{}ai-\allowbreak{}reads-\allowbreak{}human/\allowbreak{}data/\allowbreak{}energy\_\allowbreak{}raw\_\allowbreak{}m5max48/\allowbreak{}} and
\texttt{energy\_\allowbreak{}raw\_\allowbreak{}m5max128/\allowbreak{}}.}

\textbf{Result of the unisolated whole-chip attempt: net-above-idle
package subtraction could not resolve single-thread parsing on a busy
host.} On both laptops (one-minute load average 41 to 68), the pre/post
idle baseline moved by 2 to 7 W between windows while one parsing thread
added only 1 to 5 W to a chip drawing 18 to 36 W, leaving every net mean
within 1.5 standard errors of zero (several negative). That failure
motivated the dual-channel protocol in section 5c: pairing whole-chip
\texttt{IOReport} (\texttt{L2}, gated by Falsifier 6 so drifted net
subtractions are automatically rejected) with macOS kernel per-PID
\texttt{proc\_\allowbreak{}pid\_\allowbreak{}rusage(\allowbreak{}RUSAGE\_\allowbreak{}INFO\_\allowbreak{}V6).\allowbreak{}ri\_\allowbreak{}energy\_\allowbreak{}nj}
(\texttt{L3}), which attributes CPU energy strictly to the dedicated
worker subprocess executing each workload.

\section{\texorpdfstring{5c. Dual-channel Apple M5 Max
\texttt{reported\_j} across all 34 parse and tokenize
workloads}{5c. Dual-channel Apple M5 Max reported\_j across all 34 parse and tokenize workloads}}\label{c.-dual-channel-apple-m5-max-reported_j-across-all-34-parse-and-tokenize-workloads}

\emph{Executed on 2026-11-22 (\texttt{16:44:52} to
\texttt{16:53:22\ UTC}) on Apple M5 Max (\texttt{128\ GB} RAM,
\texttt{macOS\ 27.2.1} \texttt{arm64}, \texttt{Python\ 3.13.14},
\texttt{macmon\ 0.12.0}), seeded corpus \texttt{20261016}
(\texttt{20,000} records across 12 SHA-256-pinned files), \texttt{5}
interleaved shuffled rounds × \texttt{34} workloads = \texttt{170}
dedicated worker subprocess windows. Files:
\href{https://research.openie.dev/living/ai-reads-human/data/energy_reported_j.tsv}{\texttt{/\allowbreak{}living/\allowbreak{}ai-\allowbreak{}reads-\allowbreak{}human/\allowbreak{}data/\allowbreak{}energy\_\allowbreak{}reported\_\allowbreak{}j.\allowbreak{}tsv}}
and
\href{https://research.openie.dev/living/ai-reads-human/data/energy_reported_j_summary.json}{\texttt{/\allowbreak{}living/\allowbreak{}ai-\allowbreak{}reads-\allowbreak{}human/\allowbreak{}data/\allowbreak{}energy\_\allowbreak{}reported\_\allowbreak{}j\_\allowbreak{}summary.\allowbreak{}json}}.
Interactive explorer:
\href{https://research.openie.dev/living/ai-reads-human/}{\texttt{/\allowbreak{}living/\allowbreak{}ai-\allowbreak{}reads-\allowbreak{}human/\allowbreak{}}}.}

\subsection{Telemetry channels and Falsifier 6
demonstration}\label{telemetry-channels-and-falsifier-6-demonstration}

Every measurement window recorded two rootless telemetry channels
simultaneously while keeping \texttt{measured\_\allowbreak{}j\ =\ None} and
\texttt{board\_\allowbreak{}metered\ =\ false}:

\begin{enumerate}
\def\labelenumi{\arabic{enumi}.}
\tightlist
\item
  \textbf{Per-process \texttt{L3} \texttt{reported\_j} (shared-host
  label \texttt{D}, co-tenant isolated)}:
  \texttt{proc\_\allowbreak{}pid\_\allowbreak{}rusage(\allowbreak{}worker\_\allowbreak{}pid,\allowbreak{}\ RUSAGE\_\allowbreak{}INFO\_\allowbreak{}V6).\allowbreak{}ri\_\allowbreak{}energy\_\allowbreak{}nj},
  integrated by the XNU kernel strictly over the threads scheduled
  inside the dedicated worker subprocess for that single
  format/tokenizer pass.
\item
  \textbf{Whole-SoC \texttt{L2} \texttt{reported\_j} (shared-host label
  \texttt{A2}, Falsifier 6 gated)}: \texttt{macmon\ pipe} sampling Apple
  Silicon \texttt{IOReport} \texttt{chip\_power} alongside
  \texttt{0.8\ s} pre- and post-window idle baselines. Under
  \textbf{Falsifier 6}, any window whose pre/post idle baseline drifts
  by \texttt{\textgreater{}\ 5.0\%} has its
  \texttt{derived\_\allowbreak{}marginal\_\allowbreak{}j} automatically invalidated rather than
  reported. On this shared host (mean 1-minute load average
  \texttt{12.92}, idle chip draw \texttt{30.5–35.4\ W}), Falsifier 6
  invalidated \textbf{\texttt{109} of \texttt{170} windows
  (\texttt{64.1\%})} across the 34 workloads (\texttt{61\ /\ 170} valid,
  \texttt{0\ /\ 34} workloads passing all \texttt{5\ /\ 5} replicates)
  --- reproducing section 5b's diagnosis quantitatively while
  \texttt{L3} \texttt{ri\_energy\_nj} cleanly resolved every single
  workload.
\end{enumerate}

\subsection{\texorpdfstring{Parse-stage \texttt{L3} \texttt{reported\_j}
(\texttt{16} workloads, \texttt{20,000}
records)}{Parse-stage L3 reported\_j (16 workloads, 20,000 records)}}\label{parse-stage-l3-reported_j-16-workloads-20000-records}

{\def\LTcaptype{none} % do not increment counter
{\scriptsize\setlength{\tabcolsep}{3.0pt}\renewcommand{\arraystretch}{1.14}
\begin{longtable}[]{@{}
  >{\raggedright\arraybackslash}p{(\linewidth - 12\tabcolsep) * \real{0.1940}}
  >{\raggedleft\arraybackslash}p{(\linewidth - 12\tabcolsep) * \real{0.1045}}
  >{\raggedleft\arraybackslash}p{(\linewidth - 12\tabcolsep) * \real{0.1567}}
  >{\raggedleft\arraybackslash}p{(\linewidth - 12\tabcolsep) * \real{0.1567}}
  >{\raggedleft\arraybackslash}p{(\linewidth - 12\tabcolsep) * \real{0.1045}}
  >{\raggedleft\arraybackslash}p{(\linewidth - 12\tabcolsep) * \real{0.1269}}
  >{\raggedleft\arraybackslash}p{(\linewidth - 12\tabcolsep) * \real{0.1567}}@{}}
\toprule\noalign{}
\begin{minipage}[b]{\linewidth}\raggedright
\textbf{Format and parser}
\end{minipage} & \begin{minipage}[b]{\linewidth}\raggedleft
\textbf{Corpus bytes}
\end{minipage} & \begin{minipage}[b]{\linewidth}\raggedleft
\textbf{Wall ms / \texttt{20k} (mean ± SD)}
\end{minipage} & \begin{minipage}[b]{\linewidth}\raggedleft
\textbf{\texttt{L3} \texttt{nJ\ /\ record} (mean ± SD)}
\end{minipage} & \begin{minipage}[b]{\linewidth}\raggedleft
\textbf{\texttt{L3} \texttt{nJ\ /\ byte}}
\end{minipage} & \begin{minipage}[b]{\linewidth}\raggedleft
\textbf{Instructions / rec}
\end{minipage} & \begin{minipage}[b]{\linewidth}\raggedleft
\textbf{\texttt{L3} energy ratio vs binary}
\end{minipage} \\
\midrule\noalign{}
\endhead
\bottomrule\noalign{}
\endlastfoot
Fixed binary (\texttt{numpy.\allowbreak{}frombuffer}, zero-copy view) &
\texttt{460,000} & \texttt{0.\allowbreak{}0002\ ±\ 0.\allowbreak{}0000} & \texttt{0.04\ ±\ 0.01} &
\texttt{0.002} & \texttt{0.3} & \texttt{0.00×} (view only) \\
\textbf{Fixed binary (\texttt{struct.\allowbreak{}iter\_\allowbreak{}unpack}, reference)} &
\texttt{460,000} & \texttt{1.359\ ±\ 0.351} &
\textbf{\texttt{280.\allowbreak{}36\ ±\ 111.\allowbreak{}38}} & \texttt{12.19} & \texttt{1,397.4}
& \textbf{\texttt{1.00×}} \\
JSON (\texttt{orjson}, C/Rust SIMD) & \texttt{2,325,428} &
\texttt{7.053\ ±\ 1.359} & \textbf{\texttt{1,\allowbreak{}405.\allowbreak{}93\ ±\ 372.\allowbreak{}34}} &
\texttt{12.09} & \texttt{6,187.0} & \textbf{\texttt{5.01×}} \\
CSV (\texttt{csv}, Python stdlib) & \texttt{1,205,469} &
\texttt{8.804\ ±\ 1.309} & \textbf{\texttt{1,\allowbreak{}798.\allowbreak{}22\ ±\ 435.\allowbreak{}28}} &
\texttt{29.83} & \texttt{8,812.4} & \textbf{\texttt{6.41×}} \\
MessagePack (\texttt{msgpack}) & \texttt{1,660,773} &
\texttt{10.\allowbreak{}817\ ±\ 3.\allowbreak{}321} & \textbf{\texttt{1,\allowbreak{}802.\allowbreak{}46\ ±\ 593.\allowbreak{}84}} &
\texttt{21.71} & \texttt{9,593.3} & \textbf{\texttt{6.43×}} \\
\textbf{JSON (\texttt{json.loads}, Python stdlib)} & \texttt{2,325,428}
& \texttt{11.\allowbreak{}396\ ±\ 1.\allowbreak{}219} & \textbf{\texttt{2,\allowbreak{}212.\allowbreak{}12\ ±\ 409.\allowbreak{}29}} &
\texttt{19.03} & \texttt{12,700.6} & \textbf{\texttt{7.89×}} \\
JSON pretty (\texttt{json.loads}, indent 2) & \texttt{3,685,429} &
\texttt{11.\allowbreak{}716\ ±\ 2.\allowbreak{}225} & \textbf{\texttt{2,\allowbreak{}648.\allowbreak{}45\ ±\ 669.\allowbreak{}32}} &
\texttt{14.37} & \texttt{13,433.3} & \textbf{\texttt{9.45×}} \\
CBOR (\texttt{cbor2}) & \texttt{1,660,877} & \texttt{15.\allowbreak{}286\ ±\ 1.\allowbreak{}250} &
\textbf{\texttt{2,\allowbreak{}719.\allowbreak{}58\ ±\ 365.\allowbreak{}78}} & \texttt{32.75} &
\texttt{14,514.3} & \textbf{\texttt{9.70×}} \\
XML (\texttt{lxml} C library) & \texttt{2,885,487} &
\texttt{22.\allowbreak{}204\ ±\ 2.\allowbreak{}392} & \textbf{\texttt{3,\allowbreak{}800.\allowbreak{}46\ ±\ 331.\allowbreak{}49}} &
\texttt{26.34} & \texttt{19,521.4} & \textbf{\texttt{13.56×}} \\
XML (\texttt{xml.\allowbreak{}etree.\allowbreak{}ElementTree}) & \texttt{2,885,487} &
\texttt{26.\allowbreak{}354\ ±\ 0.\allowbreak{}750} & \textbf{\texttt{6,\allowbreak{}029.\allowbreak{}30\ ±\ 370.\allowbreak{}65}} &
\texttt{41.79} & \texttt{27,495.7} & \textbf{\texttt{21.51×}} \\
JSON Lines (\texttt{json.loads} per line) & \texttt{2,325,427} &
\texttt{22.\allowbreak{}729\ ±\ 1.\allowbreak{}481} & \textbf{\texttt{6,\allowbreak{}130.\allowbreak{}27\ ±\ 946.\allowbreak{}08}} &
\texttt{52.72} & \texttt{27,414.8} & \textbf{\texttt{21.87×}} \\
Markdown table (\texttt{markdown-it-py}, tokens) & \texttt{1,525,499} &
\texttt{115.\allowbreak{}160\ ±\ 7.\allowbreak{}576} & \textbf{\texttt{22,\allowbreak{}824.\allowbreak{}50\ ±\ 2,\allowbreak{}464.\allowbreak{}64}} &
\texttt{299.24} & \texttt{115,255.0} & \textbf{\texttt{81.41×}} \\
TOML (\texttt{tomllib}, Python stdlib) & \texttt{2,505,426} &
\texttt{298.\allowbreak{}582\ ±\ 60.\allowbreak{}674} & \textbf{\texttt{57,\allowbreak{}803.\allowbreak{}43\ ±\ 16,\allowbreak{}269.\allowbreak{}64}}
& \texttt{461.43} & \texttt{260,140.9} & \textbf{\texttt{206.17×}} \\
\textbf{YAML (\texttt{PyYAML} with C \texttt{libyaml})} &
\texttt{2,445,427} & \texttt{448.\allowbreak{}071\ ±\ 95.\allowbreak{}009} &
\textbf{\texttt{69,\allowbreak{}922.\allowbreak{}12\ ±\ 13,\allowbreak{}405.\allowbreak{}09}} & \texttt{571.86} &
\texttt{371,347.0} & \textbf{\texttt{249.40×}} \\
Markdown list (\texttt{markdown-it-py}, tokens) & \texttt{1,925,427} &
\texttt{417.\allowbreak{}700\ ±\ 54.\allowbreak{}289} & \textbf{\texttt{99,\allowbreak{}595.\allowbreak{}04\ ±\ 17,\allowbreak{}044.\allowbreak{}11}}
& \texttt{1,034.52} & \texttt{377,878.0} & \textbf{\texttt{355.23×}} \\
\textbf{YAML (\texttt{PyYAML}, pure Python)} & \texttt{2,445,427} &
\texttt{3,\allowbreak{}224.\allowbreak{}548\ ±\ 375.\allowbreak{}725} &
\textbf{\texttt{575,\allowbreak{}799.\allowbreak{}27\ ±\ 56,\allowbreak{}674.\allowbreak{}36}} & \texttt{4,709.19} &
\texttt{2,706,372.9} & \textbf{\texttt{2,053.75×}} \\
\end{longtable}
}
}

\subsection{\texorpdfstring{Tokenization-stage \texttt{L3}
\texttt{reported\_j} (\texttt{18} \texttt{tiktoken} workloads,
\texttt{20,000}
records)}{Tokenization-stage L3 reported\_j (18 tiktoken workloads, 20,000 records)}}\label{tokenization-stage-l3-reported_j-18-tiktoken-workloads-20000-records}

When a language model ingests text rather than structured memory,
byte-pair tokenization (\texttt{E\_tok(F,k)} in section 2) dominates CPU
decode cost on compact formats:

{\def\LTcaptype{none} % do not increment counter
{\scriptsize\setlength{\tabcolsep}{3.0pt}\renewcommand{\arraystretch}{1.14}
\begin{longtable}[]{@{}
  >{\raggedright\arraybackslash}p{(\linewidth - 12\tabcolsep) * \real{0.1884}}
  >{\raggedleft\arraybackslash}p{(\linewidth - 12\tabcolsep) * \real{0.1014}}
  >{\raggedleft\arraybackslash}p{(\linewidth - 12\tabcolsep) * \real{0.1522}}
  >{\raggedleft\arraybackslash}p{(\linewidth - 12\tabcolsep) * \real{0.1522}}
  >{\raggedleft\arraybackslash}p{(\linewidth - 12\tabcolsep) * \real{0.1014}}
  >{\raggedleft\arraybackslash}p{(\linewidth - 12\tabcolsep) * \real{0.1449}}
  >{\raggedleft\arraybackslash}p{(\linewidth - 12\tabcolsep) * \real{0.1594}}@{}}
\toprule\noalign{}
\begin{minipage}[b]{\linewidth}\raggedright
\textbf{Tokenizer \& view}
\end{minipage} & \begin{minipage}[b]{\linewidth}\raggedleft
\textbf{Tokens / \texttt{20k}}
\end{minipage} & \begin{minipage}[b]{\linewidth}\raggedleft
\textbf{Wall ms / \texttt{20k} (mean ± SD)}
\end{minipage} & \begin{minipage}[b]{\linewidth}\raggedleft
\textbf{\texttt{L3} \texttt{nJ\ /\ record} (mean ± SD)}
\end{minipage} & \begin{minipage}[b]{\linewidth}\raggedleft
\textbf{\texttt{L3} \texttt{nJ\ /\ token}}
\end{minipage} & \begin{minipage}[b]{\linewidth}\raggedleft
\textbf{Ratio vs CSV (\texttt{same\ tok})}
\end{minipage} & \begin{minipage}[b]{\linewidth}\raggedleft
\textbf{Ratio vs fixed binary parse}
\end{minipage} \\
\midrule\noalign{}
\endhead
\bottomrule\noalign{}
\endlastfoot
\texttt{cl100k\_base} · CSV (\texttt{recs.csv}) & \texttt{459,916} &
\texttt{73.\allowbreak{}442\ ±\ 20.\allowbreak{}454} & \textbf{\texttt{16,\allowbreak{}384.\allowbreak{}84\ ±\ 4,\allowbreak{}362.\allowbreak{}40}} &
\texttt{712.51} & \texttt{1.000×} & \texttt{58.44×} \\
\texttt{cl100k\_base} · Markdown table & \texttt{569,973} &
\texttt{105.\allowbreak{}364\ ±\ 18.\allowbreak{}819} & \textbf{\texttt{23,\allowbreak{}385.\allowbreak{}73\ ±\ 5,\allowbreak{}597.\allowbreak{}57}} &
\texttt{820.59} & \texttt{1.427×} & \texttt{83.41×} \\
\texttt{cl100k\_base} · Markdown list & \texttt{649,947} &
\texttt{105.\allowbreak{}199\ ±\ 11.\allowbreak{}723} & \textbf{\texttt{26,\allowbreak{}749.\allowbreak{}70\ ±\ 4,\allowbreak{}445.\allowbreak{}05}} &
\texttt{823.13} & \texttt{1.633×} & \texttt{95.41×} \\
\texttt{cl100k\_base} · JSONL (\texttt{recs.jsonl}) & \texttt{753,109} &
\texttt{154.\allowbreak{}114\ ±\ 43.\allowbreak{}130} & \textbf{\texttt{26,\allowbreak{}983.\allowbreak{}02\ ±\ 9,\allowbreak{}347.\allowbreak{}95}} &
\texttt{716.58} & \texttt{1.647×} & \texttt{96.24×} \\
\textbf{\texttt{cl100k\_base} · JSON compact (\texttt{recs.json})} &
\texttt{733,111} & \texttt{124.\allowbreak{}295\ ±\ 7.\allowbreak{}051} &
\textbf{\texttt{27,\allowbreak{}679.\allowbreak{}19\ ±\ 2,\allowbreak{}697.\allowbreak{}37}} & \texttt{755.12} &
\textbf{\texttt{1.689×}} & \textbf{\texttt{98.73×}} (\texttt{12.51×}
\texttt{json.loads}) \\
\texttt{cl100k\_base} · XML (\texttt{recs.xml}) & \texttt{1,046,596} &
\texttt{171.\allowbreak{}844\ ±\ 21.\allowbreak{}381} & \textbf{\texttt{36,\allowbreak{}730.\allowbreak{}51\ ±\ 6,\allowbreak{}968.\allowbreak{}26}} &
\texttt{701.90} & \texttt{2.242×} & \texttt{131.01×} \\
\texttt{cl100k\_base} · TOML (\texttt{recs.toml}) & \texttt{933,109} &
\texttt{178.\allowbreak{}111\ ±\ 23.\allowbreak{}760} & \textbf{\texttt{38,\allowbreak{}499.\allowbreak{}73\ ±\ 6,\allowbreak{}633.\allowbreak{}02}} &
\texttt{825.19} & \texttt{2.350×} & \texttt{137.32×} \\
\texttt{cl100k\_base} · YAML (\texttt{recs.yaml}) & \texttt{1,029,947} &
\texttt{208.\allowbreak{}434\ ±\ 33.\allowbreak{}837} & \textbf{\texttt{51,\allowbreak{}880.\allowbreak{}14\ ±\ 12,\allowbreak{}381.\allowbreak{}40}}
& \texttt{1,007.43} & \texttt{3.166×} & \texttt{185.05×} \\
\textbf{\texttt{cl100k\_base} · JSON pretty
(\texttt{recs\_\allowbreak{}pretty.\allowbreak{}json})} & \texttt{1,313,111} &
\texttt{266.\allowbreak{}827\ ±\ 21.\allowbreak{}754} & \textbf{\texttt{70,\allowbreak{}590.\allowbreak{}45\ ±\ 10,\allowbreak{}239.\allowbreak{}93}}
& \texttt{1,075.16} & \textbf{\texttt{4.308×}} &
\textbf{\texttt{251.78×}} (\texttt{26.65×} pretty
\texttt{json.loads}) \\
\texttt{o200k\_base} · CSV (\texttt{recs.csv}) & \texttt{464,995} &
\texttt{51.\allowbreak{}864\ ±\ 5.\allowbreak{}906} & \textbf{\texttt{9,\allowbreak{}570.\allowbreak{}91\ ±\ 1,\allowbreak{}012.\allowbreak{}53}} &
\texttt{411.66} & \texttt{1.000×} & \texttt{34.14×} \\
\texttt{o200k\_base} · JSON compact (\texttt{recs.json}) &
\texttt{754,794} & \texttt{93.\allowbreak{}277\ ±\ 8.\allowbreak{}521} &
\textbf{\texttt{15,\allowbreak{}125.\allowbreak{}65\ ±\ 1,\allowbreak{}287.\allowbreak{}77}} & \texttt{400.79} &
\texttt{1.580×} & \texttt{53.95×} (\texttt{6.84×}
\texttt{json.loads}) \\
\texttt{o200k\_base} · Markdown table / list & \texttt{574,818} /
\texttt{654,792} & \texttt{68.091} / \texttt{68.666} &
\textbf{\texttt{15,693.90}} / \textbf{\texttt{17,362.58}} &
\texttt{546.05} / \texttt{530.32} & \texttt{1.640×} / \texttt{1.814×} &
\texttt{55.98×} / \texttt{61.93×} \\
\texttt{o200k\_base} · JSONL / XML / TOML & \texttt{774.8k} /
\texttt{1,054.8k} / \texttt{934.8k} & \texttt{83.02} / \texttt{114.79} /
\texttt{106.68} & \textbf{\texttt{21,530.58}} /
\textbf{\texttt{22,997.48}} / \textbf{\texttt{25,307.32}} &
\texttt{555.78} / \texttt{436.05} / \texttt{541.45} &
\texttt{2.250×}--\texttt{2.644×} & \texttt{76.79×}--\texttt{90.27×} \\
\texttt{o200k\_base} · YAML (\texttt{recs.yaml}) & \texttt{1,034,792} &
\texttt{141.\allowbreak{}338\ ±\ 24.\allowbreak{}980} & \textbf{\texttt{29,\allowbreak{}401.\allowbreak{}44\ ±\ 7,\allowbreak{}308.\allowbreak{}15}} &
\texttt{568.26} & \texttt{3.072×} & \texttt{104.87×} \\
\texttt{o200k\_base} · JSON pretty (\texttt{recs\_\allowbreak{}pretty.\allowbreak{}json}) &
\texttt{1,314,794} & \texttt{193.\allowbreak{}648\ ±\ 18.\allowbreak{}210} &
\textbf{\texttt{37,\allowbreak{}070.\allowbreak{}18\ ±\ 3,\allowbreak{}651.\allowbreak{}35}} & \texttt{563.89} &
\texttt{3.873×} & \texttt{132.22×} \\
\end{longtable}
}
}

Three findings emerge from the \texttt{170}-window campaign:

\begin{enumerate}
\def\labelenumi{\arabic{enumi}.}
\tightlist
\item
  \textbf{Tokenizing JSON costs an order of magnitude more CPU
  \texttt{reported\_j} than parsing it}: \texttt{tiktoken}
  (\texttt{cl100k\_base}, implemented in compiled Rust) spends
  \texttt{27,\allowbreak{}679.\allowbreak{}2\ nJ/\allowbreak{}record} (\texttt{0.554\ J} per \texttt{20,000}
  records, \texttt{129,192} instructions/record) on compact JSON ---
  \textbf{\texttt{12.51×} Python's \texttt{json.loads()}
  (\texttt{2,\allowbreak{}212.\allowbreak{}1\ nJ/\allowbreak{}record})}, \textbf{\texttt{19.69×}
  \texttt{orjson.loads()} (\texttt{1,\allowbreak{}405.\allowbreak{}9\ nJ/\allowbreak{}record})}, and
  \textbf{\texttt{98.73×} \texttt{struct.\allowbreak{}iter\_\allowbreak{}unpack}
  (\texttt{280.\allowbreak{}4\ nJ/\allowbreak{}record})}. Indenting JSON (\texttt{indent=2})
  inflates \texttt{cl100k\_base} tokenization energy a further
  \texttt{2.55×} to \texttt{70,\allowbreak{}590.\allowbreak{}5\ nJ/\allowbreak{}record} (\texttt{4.31×} CSV
  tokenization) simply to scan whitespace runs.
\item
  \textbf{Parse energy ratios closely track instruction ratios, with
  minor effective-wattage variation (\texttt{3.28–5.65\ W} per worker
  PID)}: C \texttt{libyaml} draws \texttt{3.31\ W} effective worker
  power (\texttt{249.4×} binary \texttt{reported\_j} vs \texttt{329.8×}
  wall time), whereas per-line \texttt{jsonl} and pretty-JSON
  \texttt{tiktoken} draw \texttt{5.37–5.46\ W} (\texttt{21.87×} and
  \texttt{251.78×} \texttt{reported\_j} vs \texttt{16.73×} and
  \texttt{196.37×} wall time), confirming why direct per-process
  \texttt{ri\_energy\_nj} counters are preferable to multiplying wall
  time by a constant wattage guess.
\item
  \textbf{Provenance boundary}: \texttt{ri\_energy\_nj} is an
  operating-system/firmware counter (\texttt{reported\_j}, telemetry
  level \texttt{L3}), not an external hardware shunt or wall-plug power
  analyzer (\texttt{measured\_\allowbreak{}j\ =\ None},
  \texttt{board\_\allowbreak{}metered\ =\ false}), and still covers host CPU/memory
  work prior to neural-network forward passes
  (\texttt{e\_prefill\ ·\ T\_\{F,k\}} in M4).
\end{enumerate}

\section{6. Where a canonical machine form
fits}\label{where-a-canonical-machine-form-fits}

One response to the tax is to keep one canonical machine form for
storage and exchange and emit human text only when a person asks to read
it. Several projects have taken that route for parts of the problem:
deterministic CBOR, canonical JSON, schema-bound binary formats. The
properties such a form would need are testable: one valid byte string
per value, a declared version, closed schema, bounds on size and depth
stated before decoding, and a size claim that names the dictionary it
depends on.

\subsection{6.1 An open design question}\label{an-open-design-question}

Which canonical exchange form is best is not decided here. Canonical
JSON (RFC 8785) keeps a text form but needs rules for number formatting
and key order; deterministic CBOR (RFC 8949 section 4.2) is binary
(\texttt{2,\allowbreak{}719.\allowbreak{}6\ nJ/\allowbreak{}record} in Python \texttt{cbor2}, \texttt{9.70×}
fixed binary) but leaves choices to each protocol. A length-prefixed or
fixed-stride binary form (\texttt{280.\allowbreak{}4\ nJ/\allowbreak{}record} in
\texttt{struct.\allowbreak{}iter\_\allowbreak{}unpack}, or \texttt{0.04\ nJ/record} for a
zero-copy buffer view) removes almost all parsing work between machines,
whereas a text model still requires a projected text view
(\texttt{16,\allowbreak{}384.\allowbreak{}8\ nJ/\allowbreak{}record} for CSV vs \texttt{70,\allowbreak{}590.\allowbreak{}5\ nJ/\allowbreak{}record}
for pretty JSON under \texttt{cl100k\_base}). Which end-to-end trade-off
wins once GPU model prefill (\texttt{M4}) and field-projection tasks
(\texttt{M7}) are included remains open.

\section{7. Reproduce this page}\label{reproduce-this-page}

Everything needed is under \texttt{/\allowbreak{}living/\allowbreak{}ai-\allowbreak{}reads-\allowbreak{}human/\allowbreak{}} in this
site:

\begin{itemize}
\tightlist
\item
  \texttt{data/\allowbreak{}parse\_\allowbreak{}time.\allowbreak{}tsv}, \texttt{data/tokens.tsv},
  \texttt{data/\allowbreak{}tokenize\_\allowbreak{}time.\allowbreak{}tsv}, \texttt{data/ENV.txt},
  \texttt{data/\allowbreak{}SUMMARY.\allowbreak{}txt}: the Linux Xeon time and token baselines.
\item
  \texttt{data/\allowbreak{}energy\_\allowbreak{}reported\_\allowbreak{}j.\allowbreak{}tsv} and
  \texttt{data/\allowbreak{}energy\_\allowbreak{}reported\_\allowbreak{}j\_\allowbreak{}summary.\allowbreak{}json}: the 34-workload,
  170-window Apple M5 Max dual-channel \texttt{reported\_j} (\texttt{L3}
  \texttt{ri\_energy\_nj} + \texttt{L2} \texttt{IOReport} Falsifier 6)
  dataset (\texttt{data/\allowbreak{}energy\_\allowbreak{}joules.\allowbreak{}tsv} retains the earlier October
  2026 unisolated whole-chip null run).
\item
  \texttt{scripts/\allowbreak{}gen\_\allowbreak{}corpus.\allowbreak{}py} (seeded corpus),
  \texttt{bench\_parse.py}, \texttt{bench\_\allowbreak{}tokens.\allowbreak{}py},
  \texttt{bench\_\allowbreak{}tokenize\_\allowbreak{}time.\allowbreak{}py}, \texttt{summarize.py}: the
  measurements.
\item
  \texttt{scripts/\allowbreak{}make\_\allowbreak{}charts.\allowbreak{}py}: rebuilds every chart from the result
  files, and asserts the headline values.
\item
  \texttt{scripts/\allowbreak{}verify\_\allowbreak{}numbers.\allowbreak{}py}: recomputes every number quoted in
  the text; exits non-zero on a mismatch.
\item
  \texttt{tests-formats/}: the scripts behind the hazard test lines.
\end{itemize}

\begin{verbatim}
python3 -m venv --system-site-packages venv
venv/bin/pip install tiktoken tokenizers msgpack cbor2 orjson ruamel.yaml markdown-it-py
    psutil
venv/bin/python scripts/gen_corpus.py 20000
venv/bin/python scripts/bench_parse.py 9 > results/parse_time.tsv
venv/bin/python scripts/bench_tokens.py 2000 > results/tokens.tsv
venv/bin/python scripts/make_charts.py charts && venv/bin/python scripts/verify_numbers.py
\end{verbatim}

The tokenizer vocabularies for GPT-2, Mistral-7B-v0.1 and Qwen2.5-0.5B
are downloaded from a model hub by \texttt{bench\_\allowbreak{}tokens.\allowbreak{}py}. Timings
and \texttt{ri\_energy\_nj} counters depend on the host processor and
operating system; expect the multi-order-of-magnitude ordering across
formats to hold while absolute nanosecond and nanojoule figures shift
with hardware.

\section{8. Threats to validity}\label{threats-to-validity}

\begin{itemize}
\tightlist
\item
  Kernel \texttt{ri\_energy\_nj} (\texttt{L3} \texttt{reported\_j}) is
  an Apple Silicon model-backed operating-system counter, not an inline
  DC rail or AC wall-plug meter (\texttt{measured\_\allowbreak{}j\ =\ None},
  \texttt{board\_\allowbreak{}metered\ =\ false}), and GPU model prefill/decode
  joules (\texttt{M4}) are not yet measured here.
\item
  On shared hosts, whole-SoC \texttt{IOReport} (\texttt{L2})
  net-above-idle subtraction is vulnerable to co-tenant drift
  (\texttt{109\ /\ 170} windows invalidated by Falsifier 6 in section
  5c; null result in section 5b).
\item
  Python object-allocation overhead contributes to several parse rows
  (compiled C/Rust parsers in \texttt{M8} will bound runtime overhead;
  note \texttt{orjson} and \texttt{tiktoken} already execute compiled
  native code).
\item
  One synthetic, identifier-light corpus (\texttt{seed\ =\ 20261016});
  natural-language and telemetry logs remain untested (\texttt{M6}).
\item
  Tokenizers tested are open ones; closed commercial tokenizers are
  untested.
\item
  The 23-byte binary layout depends on a shared 24-word dictionary
  (\texttt{SHA-256} pinned in
  \texttt{data/\allowbreak{}energy\_\allowbreak{}reported\_\allowbreak{}j\_\allowbreak{}summary.\allowbreak{}json}) that per-record byte
  figures do not repeat.
\end{itemize}

\section{9. Read ledger and
references}\label{read-ledger-and-references}

Status per source is in the
\href{https://research.openie.dev/living/ai-reads-human/READ_LEDGER.md}{read
ledger} (FULL, PARTIAL, UNAVAILABLE). Cited: RFC 8259, RFC 7493, RFC
8785, RFC 8949, RFC 4180, CommonMark 0.31.2, YAML 1.2.2 (section 1.1),
TOML 1.0.0, Protocol Buffers encoding guide, OWASP CSV Injection,
\texttt{tex.web} (header only, PARTIAL), Samsi et al.~arXiv:2310.03003,
Luccioni et al.~DOI 10.1145/3630106.3658542. HTML sources were read as
tag-stripped text, so figures in them were not read.

\textbf{Not read (paywalled or unavailable), listed with identifiers and
not cited for any claim:}

{\def\LTcaptype{none} % do not increment counter
{\footnotesize\setlength{\tabcolsep}{3.8pt}\renewcommand{\arraystretch}{1.15}
\begin{longtable}[]{@{}
  >{\raggedright\arraybackslash}p{(\linewidth - 2\tabcolsep) * \real{0.5312}}
  >{\raggedright\arraybackslash}p{(\linewidth - 2\tabcolsep) * \real{0.4688}}@{}}
\toprule\noalign{}
\begin{minipage}[b]{\linewidth}\raggedright
\textbf{Item}
\end{minipage} & \begin{minipage}[b]{\linewidth}\raggedright
\textbf{Identifier}
\end{minipage} \\
\midrule\noalign{}
\endhead
\bottomrule\noalign{}
\endlastfoot
ISO 32000-1:2008, Portable Document Format & ISO catalogue
(paywalled) \\
ISO/IEC 29500, Office Open XML & ISO/IEC 29500 (free equivalent ECMA-376
only partly read) \\
ISO/IEC 26300, OpenDocument & ISO/IEC 26300 (free OASIS ODF 1.3 Part 1
read) \\
IEEE 754-2019, Floating-Point Arithmetic & DOI
10.1109/IEEESTD.2019.8766229 \\
Sennrich, Haddow, Birch, ``Neural Machine Translation of Rare Words with
Subword Units'' & DOI 10.18653/v1/P16-1162 \\
The Unicode Standard and UAX \#15 & unicode.org (not fetched) \\
ECMA-262 (number serialisation referenced by RFC 8785) & not fetched \\
The TeXbook, printed edition & ISBN 978-0-201-13447-6 \\
Horowitz, ``Computing's energy problem'', attributed source of the
quoted 640 pJ per 32-bit DRAM read & \textbf{read; 640 pJ not found in
the paper or plenary slides (they give 1 to 2 nJ, and 1.3 to 2.6 nJ per
64-bit access); not used} \\
Strubell et al.~2019; Patterson et al.~2021 and 2022; Wu et al.~2021 &
cited only inside the Samsi and Luccioni papers; not read, no claim
made \\
\end{longtable}
}
}

\section{10. What will be added}\label{what-will-be-added}

Planned expansions, in no fixed order: GPU model prefill and decode
energy per token (\texttt{M4}) and external inline wall-meter validation
(\texttt{measured\_j}); compiled C/Rust native parsers (\texttt{M8}); a
natural-language and structured-log corpus (\texttt{M6}); per-task token
counts for agent field-projection workloads (\texttt{M7}); more formats
from the ledger (PDF 2.0, C2PA, Markdown dialects, LaTeX) with their
hazards; and a worked comparison of canonical binary projections for
models.

\emph{© 2026 Open Interface Engineering, Inc.~Text, figures, data and
scripts: MIT OR Apache-2.0, at your option.}


\end{document}
